\begin{figure}[tbp]
\centering
\begin{tikzpicture}[x=.53cm,y=.53cm,font=\small]
\fill[black!5] (1.5,-.5) rectangle (2.5,8.5);
\fill[black!5] (4.5,-.5) rectangle (5.5,8.5);
\fill[black!5] (-.5,2.5) rectangle (8.5,3.5);
\fill[black!5] (-.5,5.5) rectangle (8.5,6.5);
\fill[black!12] (7.5,-.5) rectangle (8.5,8.5);
\fill[black!12] (-.5,-.5) rectangle (8.5,.5);
\foreach \i in {1,...,9}{
 \node at (\i-1,9.05) {\(\i\)};
 \node at (-1.05,9-\i) {\(\i\)};
 \foreach \j in {1,...,9}{
  \pgfmathtruncatemacro{\v}{mod(\i*\j-1,9)+1}
  \node at (\j-1,9-\i) {\(\v\)};
 }
}
\foreach \k in {0,...,9}{
 \draw[black!25] (\k-.5,-.5)--(\k-.5,8.5);
 \draw[black!25] (-.5,\k-.5)--(8.5,\k-.5);
}
\draw[line width=.65pt] (2.5,3.5) rectangle (4.5,5.5);
\node[anchor=west,align=left,text width=6.5cm] at (10,7.2) {Un soporte: \(81\) posiciones.\\[5pt]Dos evaluaciones:\\ \(S(i,j)=i+j\), \(P(i,j)=ij\).\\[5pt]Cada fila se genera por acumulación:\\ \(P(i,j+1)-P(i,j)=i\).};
\node[anchor=west,align=left,text width=6.5cm] at (10,1.8) {La carta representa \(0\pmod9\) por \(9\).\\[5pt]El bloque señalado es\\ \(B_c=\left(\begin{smallmatrix}7&2\\2&7\end{smallmatrix}\right)\),\\ con autovalores \(9\) y \(5\).};
\end{tikzpicture}
\caption{Tabla multiplicativa de APP en representantes positivos. Las filas y columnas no unitarias se distinguen por trama gris; el marco identifica el bloque diagonal adyacente de diagonales iguales. La tabla dibujada es un observable del toro aritmético; sus bordes de representación no constituyen una frontera topológica intrínseca del toro.}
\label{fig:app}
\end{figure}
